% PAGES: 39-45

\section{Lower triangular and upper triangular matrices. Tridiagonal matrices.}

\begin{definition}
A square matrix is lower triangular if and only if all the entries of the
matrix which are above the main diagonal are equal to $0$.

In other words, the $n\times n$ matrix $L$ is lower triangular if and
only if
\begin{equation}
L(j,k) \;=\; 0, \quad \forall\, 1\le j<k\le n.
\end{equation}
\end{definition}

\begin{definition}
A square matrix is upper triangular if and only if all the entries of the
matrix which are below the main diagonal are equal to $0$.

In other words, the $n\times n$ matrix $U$ is upper triangular if and
only if
\begin{equation}
U(j,k) \;=\; 0, \quad \forall\, 1\le k<j\le n.
\end{equation}
\end{definition}

From (1.103) and (1.104), it follows that the transpose of a lower
triangular matrix is an upper triangular matrix and the transpose of an
upper triangular matrix is a lower triangular matrix.

The results below follow from (1.103) and (1.104), and can be regarded as
column-based and row-based definitions for lower triangular and upper
triangular matrices:

\begin{lemma}
A square matrix of size $n$ is lower triangular if and only if the first
$k-1$ entries of the $k$-th column of the matrix are equal to $0$, for
all $k=2:n$.

In other words, the $n\times n$ matrix $L=\mathrm{col}\,(l_k)_{k=1:n}$ is
lower triangular if and only if
\begin{equation}
l_k(i) \;=\; 0, \quad \forall\, 1\le i\le k-1,\ \forall\, k=2:n.
\end{equation}

A square matrix of size $n$ is lower triangular if and only if the last
$n-j$ entries of the $j$-th row of the matrix are equal to $0$, for all
$j=1:(n-1)$.

In other words, the $n\times n$ matrix $L=\mathrm{row}\,(r_j)_{j=1:n}$ is
lower triangular if and only if
\begin{equation}
r_j(i) \;=\; 0, \quad \forall\, j+1\le i\le n,\ \forall\, j=1:(n-1).
\end{equation}
\end{lemma}

\begin{lemma}
A square matrix of size $n$ is upper triangular if and only if the last
$n-k$ entries of the $k$-th column of the matrix are equal to $0$, for
all $k=1:(n-1)$.

In other words, the $n\times n$ matrix $U=\mathrm{col}\,(u_k)_{k=1:n}$ is
upper triangular if and only if
\begin{equation}
u_k(i) \;=\; 0, \quad \forall\, k+1\le i\le n,\ \forall\, k=1:(n-1).
\end{equation}

A square matrix of size $n$ is upper triangular if and only if the first
$j-1$ entries of the $j$-th row of the matrix are equal to $0$, for all
$j=2:n$.

In other words, the $n\times n$ matrix $U=\mathrm{row}\,(r_j)_{j=1:n}$ is
upper triangular if and only if
\begin{equation}
r_j(i) \;=\; 0, \quad \forall\, 1\le i\le j-1,\ \forall\, j=2:n.
\end{equation}
\end{lemma}

Note that the sum of two lower triangular matrices is a lower triangular
matrix, the sum of two upper triangular matrices is an upper triangular
matrix, and the result of multiplying a lower triangular (or an upper
triangular) matrix by a constant is a lower triangular (or, respectively,
an upper triangular) matrix. Thus, any linear combination of lower
triangular matrices is lower triangular and any linear combination of
upper triangular matrices is upper triangular.

\begin{lemma}
(i) The product of two lower triangular matrices is lower triangular.

(ii) The product of two upper triangular matrices is upper triangular.
\end{lemma}

\begin{proof}
(i) Let $L_1=\mathrm{col}\,\!\left(l_k^{(1)}\right)_{k=1:n}$ and
$L_2=\mathrm{col}\,\!\left(l_k^{(2)}\right)_{k=1:n}$ be lower triangular
matrices. Recall from (1.105) that
\begin{align}
l_k^{(1)}(j) &= 0, \quad 1\le j\le k-1,\ \forall\, k=2:n; \\
l_k^{(2)}(j) &= 0, \quad 1\le j\le k-1,\ \forall\, k=2:n.
\end{align}

Recall from (1.11) that
\[
L_1L_2 \;=\; \mathrm{col}\,\!\left(L_1l_k^{(2)}\right)_{k=1:n}.
\]

Thus, the $k$-th column of $L_1L_2$ is $L_1l_k^{(2)}$. Using (1.7) and
the column form $L_1=\mathrm{col}\,\!\left(l_k^{(1)}\right)_{k=1:n}$ of
$L_1$, we obtain that
\begin{equation}
L_1l_k^{(2)} \;=\; \sum_{j=1}^n l_k^{(2)}(j)l_j^{(1)} \;=\; \sum_{j=k}^n l_k^{(2)}(j)l_j^{(1)},
\end{equation}
since $l_k^{(2)}(j)=0$ for all $j$ such that $1\le j\le k-1$; cf.\ (1.110).

Let $i$ such that $1\le i\le k-1$. If $k\le j\le n$, then $1\le i\le j-1$
and, from (1.109), it follows that
\begin{equation}
l_j^{(1)}(i) \;=\; 0, \quad \forall\, k\le j\le n.
\end{equation}

Then, from (1.111) and (1.112), we find that
\[
\left(L_1l_k^{(2)}\right)(i) \;=\; \sum_{j=k}^n l_k^{(2)}(j)l_j^{(1)}(i) \;=\; 0, \quad \forall\, 1\le i\le k-1.
\]

We conclude that the first $k-1$ entries of $L_1l_k^{(2)}$, which is the
$k$-th column of the matrix $L_1L_2$, are equal to $0$. Therefore, from
Lemma 1.13, it follows that $L_1L_2$ is a lower triangular matrix.

(ii) Let $U_1$ and $U_2$ be upper triangular matrices of the same size.
Recall that a matrix is upper triangular if and only if its transpose is
a lower triangular matrix. Thus, $U_1^t$ and $U_2^t$ are lower triangular
matrices. The matrix $U_2^tU_1^t$ is also a lower triangular matrix,
since we already proved that the product of two lower triangular
matrices is lower triangular. Since $(U_1U_2)^t=U_2^tU_1^t$; see (1.24),
it follows that the transpose of the matrix $U_1U_2$ is a lower
triangular matrix. We conclude that $U_1U_2$ is an upper triangular
matrix.
\end{proof}

Recall from (10.3) and (10.4), respectively, that the determinant of a
lower triangular matrix $L$ is equal to the product of all its entries on
the main diagonal, and the determinant of an upper triangular matrix $U$
is equal to the product of all its entries on the main diagonal, i.e.,
\begin{align}
\det(L) &= \prod_{i=1}^n L(i,i); \\
\det(U) &= \prod_{i=1}^n U(i,i).
\end{align}

\begin{lemma}
(i) A lower triangular matrix $L$ is nonsingular if and only if all its
entries on the main diagonal are nonzero, i.e., if and only if
\[
L(i,i) \;\ne\; 0, \quad \forall\, i=1:n.
\]

(ii) An upper triangular matrix $U$ is nonsingular if and only if all its
entries on the main diagonal are nonzero, i.e., if and only if
\[
U(i,i) \;\ne\; 0, \quad \forall\, i=1:n.
\]
\end{lemma}

\begin{proof}
These results follow from (1.113) and (1.114), and from the fact that a
matrix is nonsingular if and only if its determinant is nonzero; cf.\
Theorem 1.2.
\end{proof}

\begin{lemma}
(i) The inverse of an upper triangular matrix is upper triangular.

(ii) The inverse of a lower triangular matrix is lower triangular.
\end{lemma}

A proof of this result can be found in Section 10.7; see Lemma 10.18.

\begin{definition}
The $n\times n$ matrix $A$ is banded of band $m$ if and only if
\[
A(j,k) \;=\; 0, \quad \forall\, 1\le j,k\le n \text{ with } |j-k|>m.
\]
\end{definition}

Note that a banded matrix of band $0$ is a diagonal matrix, since
$A(j,k)=0$ for all $j$ and $k$ with $|j-k|>0$ is equivalent to $A(j,k)=0$
for all $j\ne k$.

\begin{examplex}
The positions of the entries of a $6\times 6$ banded matrix with band $2$
which could be nonzero are denoted by ``$\times$'' in the matrix below:
\[
\begin{pmatrix}
\times & \times & \times & 0 & 0 & 0 \\
\times & \times & \times & \times & 0 & 0 \\
\times & \times & \times & \times & \times & 0 \\
0 & \times & \times & \times & \times & \times \\
0 & 0 & \times & \times & \times & \times \\
0 & 0 & 0 & \times & \times & \times
\end{pmatrix}.
\]
\end{examplex}
